\section{Discrete sale choices and continuous financing}
The inventory offers give three sale choices: none, half or all. They do not
make the financing policy set finite. A principal cap permits an interval of
bridge draws unless the contract also specifies denominations. To compare
policies, retain those amounts as variables and first separate the cases in
which each offer is used. For each feasible policy, record minimum cash, final
cash, fees, remaining debt and inventory. A policy minimizing fees may differ
from one preserving inventory or maximizing the lowest cash balance.

The action domain can be derived directly for this example. Let $z$ be the
fraction sold, with $z\in\{0,1/2,1\}$, and let $a$ be bridge principal. Every
sale choice still needs the bridge before day 4. Write $b$ for term face and
$u$ for whether B is issued: $u=0$ requires $b=0$, while $u=1$ requires
$0<b\leq40$. Hold the supplied price $q>0$ fixed. For this calculation, both
sale and bridge proceeds arrive before the old debt payment; any day-5 customer
receipt arrives before B settles, and B's fee is withheld from its proceeds.
These are explicit settlement-order assumptions. They do not give the manager
permission to decide an issue after its contractual notice deadline.

Let $R$ be the cash collected by day 6: $R=18$ in H1 and $R=0$ in H2.
Cash just after the old debt payment is
\begin{equation}
 C_4=12-10+15z+(a-1)-32=a+15z-31.
\end{equation}
Between that payment and bridge repayment, the company receives $R$ and the
term issue supplies net cash $qb-2u$. Repaying principal plus one therefore gives
\begin{equation}
 C_6=C_4+R+(qb-2u)-(a+1)
     =15z+R+qb-32-2u.
\end{equation}
The bridge principal cancels from $C_6$ because the company repays the extra
amount it borrowed. It still changes $C_4$. Under the stated event order,
nonnegative cash throughout days 1--6 is equivalent to
\begin{equation}
 31-15z\leq a\leq30,
 \qquad 15z+R+qb-32-2u\geq0,
\end{equation}
together with the issue/no-issue domain and offer eligibility. To see why no
other cash constraint is missing, cash after wages is 2 and the sale and bridge
then add nonnegative proceeds. After day 4, the receipt adds cash; cash after B
settles must be $C_6+a+1$, which is positive whenever $C_6\geq0$. The bridge
repayment is the final outflow in this comparison. H2's day-9 receipt increases
cash later; it cannot repair a negative prefix on day 6. Any term debt remains
an obligation beyond this cash horizon.

With the inventory covenant, $u=1$ also requires $20(1-z)\geq10$.
For half sale, the bridge interval is $[23.5,30]$. If B is used, its face must
satisfy $qb\geq8.5$ in H1 or $qb\geq26.5$ in H2, recovering the earlier price
thresholds at the face cap. No sale would require $a\geq31$, so that entire
branch is infeasible. Full sale permits $a\in[16,30]$; with the covenant it
must use $u=0$, giving $C_6=1$ in H1 and $C_6=-17$ in H2. Without the covenant,
the term issue remains a separate possible response.

These intervals describe the continuous amounts; a few selected amounts are
only witnesses. For a fixed sale choice and term issue, taking the least
bridge principal $a=31-15z$ preserves feasibility under these assumptions, but
it can change the best liquidity policy. In H1 with full sale and no term issue,
$a=16$ leaves cash 0 after the old debt payment and cash 1 after bridge
repayment. Taking $a=17$ instead leaves cash 1 after both payments. Both paths
pay the same two units of bridge cost and end with the same cash and no bridge
debt, yet their minimum cash balances are 0 and 1. The latter is maximal within
this full-sale/no-term comparison because every path ends with cash 1. Restricting
the comparison to minimum funding would lose that conclusion. An interest
charge depending on the amount borrowed would require a different calculation.

Settlement order also matters for the domain itself. Without the inventory
covenant, take full sale, H1,
$a=16$, $q=0.9$ and $b=10/9$. Issuing B yields net cash $-1$ after its fixed
fee. Receipt before issue gives a feasible path ending at cash 0. Issue before
receipt instead asks the company to pay that one unit while cash is still 0.
That order requires the additional prefix condition $C_4+qb-2\geq0$.
The small issue is not proposed as an attractive policy; it exposes an action
that a daily net total would mistakenly admit under the opposite order.

A finite enumeration is therefore complete only for an explicitly finite
action menu or after a justified reduction for the stated objective. Here one
can enumerate sale choices and offer activation, then retain or solve the
continuous inequalities inside each branch. With fixed $q$ these are linear
constraints. If price depends on face, inventory or the later policy, that
continuous problem changes. Across H1 and H2, decisions must still share the
same values until the manager's observations permit them to differ. Comparison
with another optimizer requires the same amounts, information, settlement
order and objective; a better result under different choices answers a
different financing question.
